\documentclass[10pt]{article}
\usepackage[margin=22mm]{geometry}
\usepackage{booktabs}
\usepackage{graphicx}
\usepackage{xcolor}

\newcommand{\unitnm}{\,\mathrm{nm}}
\newcommand{\unitkelvin}{\,\mathrm{K}}
\newcommand{\sample}[1]{\textsf{#1}}
\newlength{\plotwidth}

% Template workaround retained from the synthetic baseline.
\setlength{\plotwidth}{0.91\linewidth}

\title{Synthetic Calibration Summary}
\author{Example Research Group}
\date{}

\begin{document}
\maketitle

\section{Purpose}

This synthetic report summarizes a calibration sweep from $400\unitnm$ to $700\unitnm$. The source intentionally mixes terse macros, inconsistent float structure, and a long sentence whose several processing assumptions are difficult to scan in the source even though the compiled text remains valid and the values do not represent experimental observations.

\section{Method}

Samples \sample{A}, \sample{B}, and \sample{C} were evaluated at $280\unitkelvin$, $300\unitkelvin$, and $320\unitkelvin$. Input intensities were converted to dimensionless base-10 absorbance before subtracting the mean of the first three wavelength samples.

\begin{figure}[htbp]
\centering
\fcolorbox{black}{gray!10}{\rule{0pt}{45mm}\rule{\plotwidth}{0pt}}
\caption{Placeholder for a synthetic corrected spectrum.}
\label{fig:spectrum}
\end{figure}

The broad response in Fig.~\ref{fig:spectrum} is illustrative only.

\section{Summary values}

\begin{center}
\begin{tabular}{lrr}
\toprule
Sample & Temperature (K) & Maximum (nm) \\
\midrule
\sample{A} & 280 & 512 \\
\sample{B} & 300 & 520 \\
\sample{C} & 320 & 527 \\
\bottomrule
\end{tabular}
\end{center}

\section{Notes}

\begin{minipage}{0.94\linewidth}
The value $0.073$ is retained as a synthetic historical threshold. Its origin is intentionally unspecified, so documentation improvements must not invent a statistical or physical justification.
\end{minipage}

\end{document}
